\documentclass[11pt]{letter}

\usepackage[margin=1in]{geometry}
\usepackage[
  colorlinks=true,
  linkcolor=blue,
  citecolor=blue,
  urlcolor=blue
]{hyperref}

\signature{Celia Blanco\\On behalf of all authors}
\address{Celia Blanco\\Centro de Astrobiolog\'{i}a (CSIC-INTA)\\Ctra de Torrej\'{o}n a Ajalvir, km 4\\28850 Torrej\'{o}n de Ardoz, Madrid, Spain\\\href{mailto:celia.blanco@cab.inta-csic.es}{celia.blanco@cab.inta-csic.es}}

\begin{document}

\begin{letter}{Technological Forecasting \& Social Change}

\opening{Dear Editor,}

We are pleased to submit our manuscript entitled ``Projections of Earth's Technosphere: Civilization Collapse--Recovery Dynamics and Detectability'' for consideration in \textit{Technological Forecasting \& Social Change}.

This paper is part of the ``Projections of Earth's Technosphere'' series. It builds on the scenario modeling and technosignature framework established by Haqq-Misra et al.\ (\textit{Technol.\ Forecast.\ Soc.\ Change}, 2025; \href{https://doi.org/10.1016/j.techfore.2025.124194}{doi:10.1016/j.techfore.2025.124194}) and introduces a quantitative simulation layer that models collapse--recovery dynamics across the ten scenarios defined in that earlier work. Using a hybrid deterministic-stochastic framework over a 1000-year window, we characterize the ``duty cycle'' of technological civilizations---the fraction of their lifespan during which they remain technologically active---and identify the parameters that most strongly govern long-term trajectories.

The manuscript serves a dual purpose. First, it identifies actionable resilience levers for Earth's civilization: sensitivity analysis reveals that reducing resource depletion and preserving post-collapse recovery capacity are, in most scenarios, the most impactful resilience levers. Second, it contributes to the search for extraterrestrial technosignatures by deriving an effective detectability duration that accounts for intermittent civilizational activity, offering a quantitative explanation for the apparent absence of extraterrestrial signals that does not require intelligent life to be rare.

We believe this work is well suited to \textit{Technological Forecasting \& Social Change} given its integration of futures scenario analysis, civilizational resilience modeling, and long-term technological forecasting.

The manuscript has not been published elsewhere and is not under consideration by any other journal. All authors have approved the manuscript and agree with its submission.

\closing{Sincerely,}

\end{letter}
\end{document}
